\small
\begin{enumerate}[label=\textbf{\arabic*.}, leftmargin=0pt, itemsep=1.2em]
\item Как обычно представляют две вещественные переменные в бинарном ГА для двумерной оптимизации?
\begin{itemize}[label=$\square$, leftmargin=1.5cm, itemsep=0.2em, topsep=0.3em]
\item Как конкатенацию двух бинарных подстрок одинаковой длины
\item Как одну вещественную матрицу без кодирования
\item Как набор меток классов
\item Как последовательность только из нулей
\end{itemize}
{\small\textit{Правильный ответ: Как конкатенацию двух бинарных подстрок одинаковой длины.}}

\item Что означает декодирование бинарной хромосомы для двух переменных?
\begin{itemize}[label=$\square$, leftmargin=1.5cm, itemsep=0.2em, topsep=0.3em]
\item Преобразование двух битовых подстрок в значения x1 и x2 заданного диапазона
\item Удаление половины битов хромосомы
\item Перестановку местами двух переменных
\item Переход от евклидовой метрики к манхэттенской
\end{itemize}
{\small\textit{Правильный ответ: Преобразование двух битовых подстрок в значения x1 и x2 заданного диапазона.}}

\item В чем идея однородного кроссинговера?
\begin{itemize}[label=$\square$, leftmargin=1.5cm, itemsep=0.2em, topsep=0.3em]
\item Каждый бит потомка выбирается от одного из родителей согласно случайной маске
\item Выбирается только одна точка разреза
\item Реверсируется случайный сегмент хромосомы
\item Родители просто копируются без изменений
\end{itemize}
{\small\textit{Правильный ответ: Каждый бит потомка выбирается от одного из родителей согласно случайной маске.}}

\item Что делает инверсионная мутация в бинарной хромосоме?
\begin{itemize}[label=$\square$, leftmargin=1.5cm, itemsep=0.2em, topsep=0.3em]
\item Разворачивает выбранный сегмент битов в обратном порядке
\item Заменяет все нули на единицы во всей строке
\item Меняет местами две популяции
\item Удаляет худшую особь
\end{itemize}
{\small\textit{Правильный ответ: Разворачивает выбранный сегмент битов в обратном порядке.}}

\item Какое свойство сохраняется при инверсионной мутации?
\begin{itemize}[label=$\square$, leftmargin=1.5cm, itemsep=0.2em, topsep=0.3em]
\item Общее число битов и набор битов в хромосоме
\item Значение фитнеса особи
\item Координаты глобального минимума
\item Размер турнира
\end{itemize}
{\small\textit{Правильный ответ: Общее число битов и набор битов в хромосоме.}}

\item Как работает турнирная селекция?
\begin{itemize}[label=$\square$, leftmargin=1.5cm, itemsep=0.2em, topsep=0.3em]
\item Из случайно выбранной группы особей выбирается лучшая
\item Родитель всегда выбирается случайно без учета качества
\item Из популяции удаляется лучшая особь
\item Потомки генерируются только из одной особи
\end{itemize}
{\small\textit{Правильный ответ: Из случайно выбранной группы особей выбирается лучшая.}}

\item Для чего нужна стратегия сокращения промежуточной популяции?
\begin{itemize}[label=$\square$, leftmargin=1.5cm, itemsep=0.2em, topsep=0.3em]
\item Чтобы выбрать, какие особи останутся в следующем поколении
\item Чтобы уменьшить длину каждой хромосомы
\item Чтобы вычислить псевдообратную матрицу
\item Чтобы нормализовать входные данные
\end{itemize}
{\small\textit{Правильный ответ: Чтобы выбрать, какие особи останутся в следующем поколении.}}

\item В чем особенность элитарной стратегии сокращения?
\begin{itemize}[label=$\square$, leftmargin=1.5cm, itemsep=0.2em, topsep=0.3em]
\item Лучшие особи гарантированно сохраняются
\item Всегда сохраняются только потомки
\item Сохраняются только случайные особи
\item Удаляются все родители и потомки одновременно
\end{itemize}
{\small\textit{Правильный ответ: Лучшие особи гарантированно сохраняются.}}

\item Что характерно для стратегии полной замены?
\begin{itemize}[label=$\square$, leftmargin=1.5cm, itemsep=0.2em, topsep=0.3em]
\item Новое поколение формируется исключительно из потомков
\item Всегда сохраняется один родитель
\item Потомки никогда не проходят отбор
\item Сохраняются только худшие особи
\end{itemize}
{\small\textit{Правильный ответ: Новое поколение формируется исключительно из потомков.}}

\item Что характерно для случайной замены?
\begin{itemize}[label=$\square$, leftmargin=1.5cm, itemsep=0.2em, topsep=0.3em]
\item Часть родителей вытесняется потомками случайным образом
\item Сохраняются только лучшие родители
\item Ни одна особь не может быть удалена
\item Кроссинговер запрещается
\end{itemize}
{\small\textit{Правильный ответ: Часть родителей вытесняется потомками случайным образом.}}

\item Почему функция Шаффера F7 считается сложной для оптимизации?
\begin{itemize}[label=$\square$, leftmargin=1.5cm, itemsep=0.2em, topsep=0.3em]
\item Она имеет много локальных минимумов
\item Она линейна на всей области
\item У нее отсутствует глобальный минимум
\item Она определена только для целых значений
\end{itemize}
{\small\textit{Правильный ответ: Она имеет много локальных минимумов.}}

\item Где расположен глобальный минимум функции Шаффера F7 в рассматриваемой постановке?
\begin{itemize}[label=$\square$, leftmargin=1.5cm, itemsep=0.2em, topsep=0.3em]
\item В точке (0, 0)
\item В точке (1, 1)
\item На границе области
\item Он неизвестен даже теоретически
\end{itemize}
{\small\textit{Правильный ответ: В точке (0, 0).}}

\item Какое преимущество дает бинарное кодирование двух переменных?
\begin{itemize}[label=$\square$, leftmargin=1.5cm, itemsep=0.2em, topsep=0.3em]
\item Позволяет использовать классические битовые операторы кроссинговера и мутации
\item Исключает необходимость декодирования
\item Делает функцию всегда выпуклой
\item Устраняет локальные минимумы
\end{itemize}
{\small\textit{Правильный ответ: Позволяет использовать классические битовые операторы кроссинговера и мутации.}}

\item Что определяет разрешающую способность бинарного представления по каждой переменной?
\begin{itemize}[label=$\square$, leftmargin=1.5cm, itemsep=0.2em, topsep=0.3em]
\item Число бит, выделенных под переменную
\item Только размер популяции
\item Только вероятность кроссинговера
\item Только число поколений
\end{itemize}
{\small\textit{Правильный ответ: Число бит, выделенных под переменную.}}

\item Зачем генетическому алгоритму в мультимодальной задаче нужно разнообразие популяции?
\begin{itemize}[label=$\square$, leftmargin=1.5cm, itemsep=0.2em, topsep=0.3em]
\item Чтобы не застрять в локальных экстремумах
\item Чтобы ускорить вычисление евклидова расстояния
\item Чтобы убрать необходимость в селекции
\item Чтобы сделать целевую функцию линейной
\end{itemize}
{\small\textit{Правильный ответ: Чтобы не застрять в локальных экстремумах.}}

\item Какой эффект обычно дает слишком сильное давление отбора?
\begin{itemize}[label=$\square$, leftmargin=1.5cm, itemsep=0.2em, topsep=0.3em]
\item Потерю разнообразия и преждевременную сходимость
\item Гарантированный выход к глобальному минимуму
\item Отсутствие потомков
\item Переход к обучению без учителя
\end{itemize}
{\small\textit{Правильный ответ: Потерю разнообразия и преждевременную сходимость.}}

\item Для чего обычно анализируют разные стратегии редукции популяции?
\begin{itemize}[label=$\square$, leftmargin=1.5cm, itemsep=0.2em, topsep=0.3em]
\item Чтобы оценить их влияние на сходимость и устойчивость поиска
\item Чтобы выбрать функцию активации нейрона
\item Чтобы определить число входных признаков
\item Чтобы вычислить коэффициент силуэта
\end{itemize}
{\small\textit{Правильный ответ: Чтобы оценить их влияние на сходимость и устойчивость поиска.}}

\item Какой общий смысл имеет функция приспособленности в задаче минимизации?
\begin{itemize}[label=$\square$, leftmargin=1.5cm, itemsep=0.2em, topsep=0.3em]
\item Чем лучше особь, тем меньше значение целевой функции или тем выше специально определенный фитнес
\item Она всегда должна быть равна длине хромосомы
\item Она не зависит от x1 и x2
\item Она вычисляется только после завершения алгоритма
\end{itemize}
{\small\textit{Правильный ответ: Чем лучше особь, тем меньше значение целевой функции или тем выше специально определенный фитнес.}}

\item Что позволяет визуализация популяции на линиях уровня функции?
\begin{itemize}[label=$\square$, leftmargin=1.5cm, itemsep=0.2em, topsep=0.3em]
\item Наблюдать, как решения перемещаются к перспективным областям
\item Доказать линейную разделимость данных
\item Измерить скорость работы процессора
\item Вычислить точное значение градиента
\end{itemize}
{\small\textit{Правильный ответ: Наблюдать, как решения перемещаются к перспективным областям.}}

\item Как обычно влияет увеличение числа поколений?
\begin{itemize}[label=$\square$, leftmargin=1.5cm, itemsep=0.2em, topsep=0.3em]
\item Дает алгоритму больше возможностей улучшить решение
\item Сразу уменьшает размер хромосомы
\item Отменяет необходимость мутации
\item Гарантирует отсутствие локальных минимумов
\end{itemize}
{\small\textit{Правильный ответ: Дает алгоритму больше возможностей улучшить решение.}}

\item Что дает увеличение размера популяции в двумерной задаче оптимизации?
\begin{itemize}[label=$\square$, leftmargin=1.5cm, itemsep=0.2em, topsep=0.3em]
\item Лучшее покрытие пространства поиска
\item Уменьшение числа генов в хромосоме
\item Отказ от кроссинговера
\item Переход от бинарного к вещественному кодированию
\end{itemize}
{\small\textit{Правильный ответ: Лучшее покрытие пространства поиска.}}

\item Почему в бинарном ГА важно корректно декодировать обе переменные отдельно?
\begin{itemize}[label=$\square$, leftmargin=1.5cm, itemsep=0.2em, topsep=0.3em]
\item Иначе невозможно восстановить точку пространства поиска
\item Иначе мутация перестает работать математически
\item Иначе пропадает элитарная стратегия
\item Иначе турнирная селекция становится случайной
\end{itemize}
{\small\textit{Правильный ответ: Иначе невозможно восстановить точку пространства поиска.}}

\item Какой из операторов непосредственно комбинирует информацию от двух родителей?
\begin{itemize}[label=$\square$, leftmargin=1.5cm, itemsep=0.2em, topsep=0.3em]
\item Кроссинговер
\item Мутация
\item Редукция
\item Оценка фитнеса
\end{itemize}
{\small\textit{Правильный ответ: Кроссинговер.}}

\item Какой из операторов чаще всего вводит новые структуры, отсутствовавшие у родителей?
\begin{itemize}[label=$\square$, leftmargin=1.5cm, itemsep=0.2em, topsep=0.3em]
\item Мутация
\item Элитарный отбор
\item Ранжирование популяции
\item Декодирование
\end{itemize}
{\small\textit{Правильный ответ: Мутация.}}

\item Что можно сказать о применимости битовых операторов к бинарному кодированию двух переменных?
\begin{itemize}[label=$\square$, leftmargin=1.5cm, itemsep=0.2em, topsep=0.3em]
\item Они естественно работают с кодом независимо от смысла отдельных переменных
\item Они работают только для одной переменной
\item Они запрещены в мультимодальных задачах
\item Они требуют дифференцируемой целевой функции
\end{itemize}
{\small\textit{Правильный ответ: Они естественно работают с кодом независимо от смысла отдельных переменных.}}

\item В чем состоит компромисс между исследованием и эксплуатацией в ГА?
\begin{itemize}[label=$\square$, leftmargin=1.5cm, itemsep=0.2em, topsep=0.3em]
\item Нужно сочетать сохранение хороших решений и поиск новых областей
\item Нужно одновременно минимизировать и максимизировать одну и ту же функцию
\item Нужно полностью исключить отбор
\item Нужно всегда выбирать максимально возможную мутацию
\end{itemize}
{\small\textit{Правильный ответ: Нужно сочетать сохранение хороших решений и поиск новых областей.}}

\item Какой оператор особенно полезен для перестройки внутренней структуры бинарной строки без добавления новых битов?
\begin{itemize}[label=$\square$, leftmargin=1.5cm, itemsep=0.2em, topsep=0.3em]
\item Инверсия сегмента
\item Рулеточная селекция
\item Элитизм
\item Декодирование
\end{itemize}
{\small\textit{Правильный ответ: Инверсия сегмента.}}

\item Какой результат в мультимодальной задаче обычно считается успешным?
\begin{itemize}[label=$\square$, leftmargin=1.5cm, itemsep=0.2em, topsep=0.3em]
\item Нахождение области, близкой к глобальному минимуму
\item Обязательное совпадение со всеми локальными минимумами
\item Получение нулевой длины хромосомы
\item Отсутствие мутаций в процессе поиска
\end{itemize}
{\small\textit{Правильный ответ: Нахождение области, близкой к глобальному минимуму.}}

\item Почему параметры ГА подбирают экспериментально?
\begin{itemize}[label=$\square$, leftmargin=1.5cm, itemsep=0.2em, topsep=0.3em]
\item Потому что эффективность зависит от конкретной задачи и формы пространства поиска
\item Потому что все параметры математически фиксированы
\item Потому что мутация всегда должна быть 0.5
\item Потому что бинарное кодирование всегда дает оптимум без настройки
\end{itemize}
{\small\textit{Правильный ответ: Потому что эффективность зависит от конкретной задачи и формы пространства поиска.}}

\item Какое утверждение о бинарном ГА для двух переменных наиболее корректно?
\begin{itemize}[label=$\square$, leftmargin=1.5cm, itemsep=0.2em, topsep=0.3em]
\item Он подходит для поиска в сложных негладких пространствах, не требуя вычисления производных
\item Он работает только для выпуклых функций
\item Он требует наличия обучающей выборки с метками
\item Он не использует случайность
\end{itemize}
{\small\textit{Правильный ответ: Он подходит для поиска в сложных негладких пространствах, не требуя вычисления производных.}}

\end{enumerate}
